\section{Existence}\label{sec:existence}

We give two constructions that turn hypotheses on the data into a realization. Neither assumes that a perturbed
surface embeds as a whole. Theorem~\ref{thm:lemma0} does assume (P0) that each original patch is embedded, and two
of its hypotheses are non-local.

\subsection{Fans and cones}\label{sec:fans}

\subsubsection*{Setting}
$K$ is regular (every face a disk bounded by a single wire). For B-rep complexes:
\begin{itemize}
\item at a seam edge, E3$'$ below refers to the two sheets of the face;
\item at a closed edge, V1 counts edge ends;
\item for a face with inner wires, P0 becomes ``$Q_f$ is homeomorphic to the face''.
\end{itemize}
$\mathrm{dev}^*_e(t)$ is the deviation of the boundary points $b_{f,e}(t)=\partial P_f\cap\Pi(t)$ from $\gamma_e(t)$ in
the normal plane $\Pi(t)$; it is at most $\mathrm{dev}_e$.

\subsubsection*{Auxiliary sets}
\begin{itemize}
\item \emph{Normal disks.} For an edge $e$, choose parameters $0<a_e<b_e<1$ and a $C^1$ radius function $\rho_e(t)$.
$\rho_e$ is constant on the stubs $[t_v,a_e]$, $[b_e,t_w]$ and on a neighbourhood of $a_e$ and $b_e$; here $t_v$ is
the last parameter with $\gamma_e(t)\in B(p_v,r_v)$. $D_e(t)$ is the closed disk of radius $\rho_e(t)$ in $\Pi(t)$
centred at $\gamma_e(t)$, and $T_e=\bigcup_{t\in[a_e,b_e]}D_e(t)$.
\item \emph{Vertex regions.} For a vertex $v$, choose $r_v$. The vertex region $W_v$ is the union of the ball
$B(p_v,r_v)$ and, for each edge $e$ at $v$, the disks $D_e(t)$ for $t$ between $t_v$ and $a_e$ (or $b_e$).
\end{itemize}

\subsubsection*{Hypotheses}
E1--E3$'$ are required on $[t_v,t_w]$, the stubs included.
\begin{description}
\item[E1] $\mathrm{dev}^*_e(t)<\rho_e(t)<\min(\tau_e,\operatorname{reach}\gamma_e|_{[a_e,b_e]})$.
\item[E2] For every face $f>e$:
  \begin{enumerate}
  \item $\psi_f<\pi/2$, where $\psi_f$ is the largest angle between $\gamma_e'$ and $T P_f$ inside the tube;
  and $\rho_e k_f<1$, where $k_f=\kappa_f/\cos\psi_f$;
  \item $t\mapsto\pi_e(\beta_{f,e}(t))$ is strictly monotone. A sufficient condition is $L_e+\kappa_\gamma
  \mathrm{dev}^*_e<1$, where $L_e$ is the arc-length Lipschitz constant of $\beta_{f,e}-\gamma_e$;
  \item \emph{no return:} $\dist(\gamma_e(t),P_f\setminus C_f)>\rho_e(t)$. Here $C_f$ is the union of the section
  arcs of $P_f$ from $b_{f,e}$ to their exit points, of length at most
  $\ell_f=(\mathrm{dev}^*+\sqrt{\mathrm{dev}^{*2}+\mu\rho^2})/\mu$ with $\mu=1-\rho k_f$.
  \end{enumerate}
\item[E3$'$] The exit points $k_{f,e}(t)$ and $k_{g,e}(t)$ of the two face sections on $\partial D_e(t)$ are distinct.
A sufficient condition (E3) is $\alpha_e(t)>\varepsilon_f+\varepsilon_g$, with
$\varepsilon=\arcsin\min\{1,(\mathrm{dev}^*_e+k\ell^2/2)/\rho_e\}$. Here $\alpha_e$ is the angle between the inward
tangents of the two sections.
\item[V0] $\gamma_e$'s end at $v$ and the corners of the patches at $v$ lie in $\inter B(p_v,r_v)$.
\item[V1] (a) $W_v$ is strictly star-shaped with respect to $p_v$, and $W_v\subseteq\inter\Z(v)$.
(b) \emph{(non-local)} The construction outside $W_v$ meets $\partial W_v$ transversally in one simple closed curve
$\lambda_v$. Each incident edge contributes one point and each incident face one arc, in the cyclic order of the
link of $v$.
\item[S1] (a) The $W_v$ are pairwise disjoint. (b) The $T_e$ are pairwise disjoint. (c) $\gamma_e$ and $T_e$ miss $W_w$
for $w\notin\partial e$. (d) $P_f$ misses $T_e$ for $e\not<f$, and misses $W_v$ for $v\not<f$. (e) Outside
$\bigcup T\cup\bigcup W$, the patch pieces are pairwise disjoint. (f) $T_e\cap W_v=D_e(a_e)$ for $v<e$, and each
piece outside $W_v$ meets $W_v$ only in $\lambda_v$.
\item[P0] \emph{(non-local)} Each $P_f$ is embedded, and $P_f\cap(\bigcup T\cup\bigcup W)$ is a collar of $\partial
P_f$, so that $Q_f=P_f\setminus\inter(\bigcup T\cup\bigcup W)$ is a closed disk.
\end{description}
V1(b) and P0 are not consequences of local data. In the models of Section~\ref{sec:experiments} they were verified
only for the gapped cube, implicitly, through an exact check that the constructed surface is embedded. The
sampled checker does not test V1(b), P0, E2.3 beyond S1(d), or S1(e) for adjacent faces.

\begin{theorem}[Existence by fans and cones; conditional, proof outline]\label{thm:lemma0}
Under E1--E3$'$, V0, V1, S1 and P0, the tolerant B-rep has a realization.
\end{theorem}
\begin{proof}[Construction and proof outline]
\begin{enumerate}
\item \emph{Vertices and edges.} Put $h(v)=p_v$. On $[a_e,b_e]$, let $h(e)$ follow $\gamma_e$.
\item \emph{Fans.} In each disk $D_e(t)$, the radius from $\gamma_e(t)$ to $k_{f,e}(t)$ is assigned to $f$; these
radii form the fan strip $F_{f,e}$.
\item \emph{Faces away from the vertices.} Outside the vertex regions, the face is $Q_f\cup\bigcup_eF_{f,e}$.
\item \emph{Vertex regions.} Inside $W_v$, take the cone from $p_v$ over $\lambda_v$.
\item \emph{Injectivity.}
  \begin{itemize}
  \item The disks $D_e(t)$ are pairwise disjoint (E1 and the reach), and within one disk there are two radii meeting
  only at the centre (E3$'$).
  \item Nothing else enters the tubes or vertex regions (S1, P0, and E2.3 for $Q_f\cap F_{f,e}=k_{f,e}$).
  \item The cone is injective by star-shapedness and V1(b).
  \end{itemize}
\item \emph{Cell structure and zones.} The pieces of a face glue to a closed disk with boundary $h(\partial f)$, in
the order forced by P0 and V1(b). Each piece lies in the zone of its cell with a positive margin.
\item \emph{PL step.} The map is bi-Lipschitz: every crease has a positive opening angle, and at points where three
or more pieces meet, the tangent cone is an embedded disk. A fine secant approximation on a shape-regular
triangulation is then injective and stays in the zones.
\end{enumerate}
\draftnote{The full argument for steps 5--7 is in the repository (\texttt{notes/lemma0\_v02.md} \S\S4--5, with the
changes of \texttt{lemma0\_v03.md} \S1); it should be written out here. One gap is open: the apex estimate in the PL
step quotes a curvature $\kappa_\lambda$ of $\lambda_v$, but $\lambda_v$ is only $C^1$.}
\end{proof}

\begin{example}
On a cube whose faces are pulled apart by gaps, the construction was carried out exactly in rational
arithmetic: 216 triangles, with zone membership and embeddedness checked exactly
(\texttt{compute/lemma0\_box.py}).
\end{example}

\subsection{Re-intersection}

\begin{proposition}[Planar, degree three]\label{prop:B}
Suppose every face is a disk bounded by a single wire and is planar, $P_f\subseteq\Pi_f$, and every vertex lies on exactly three faces whose planes are
independent. Let $c_v$ be the common point of the three planes, $\ell_e=[c_v,c_w]$ for an edge $e=vw$, and
$\tilde P_f$ the polygon on the $c$'s of the vertices of $f$, in the order of $\partial f$. Assume:
\begin{itemize}
\item $c_v\in\Z(v)$, $\ell_e\subseteq\Z(e)$ and $\tilde P_f\subseteq\Z(f)$;
\item the $\tilde P_f$ are simple polygons that meet exactly in their common edges and vertices.
\end{itemize}
Then the cellular map onto $\bigcup\tilde P_f$ is a realization.
\end{proposition}
\begin{proof}
It is a PL homeomorphism onto $\bigcup\tilde P_f$, and each cell lands in its zone.
\end{proof}

\begin{proposition}[Curved, transversal]\label{prop:Bprime}
Suppose every face lies on a $C^2$ surface $S_f$, and there is a \emph{re-trimmed model}: patches $\tilde P_f\subseteq
S_f$, curves $\tilde\gamma_e$ and points $c_v$ bounding a solid, whose cell structure is isomorphic to $K$ (as a
B-rep complex, respecting seams and closed edges). Assume:
\begin{itemize}
\item the closed cells lie in their zones;
\item the cells are $C^2$ up to the boundary, the curves are regular, and there are no cusp corners;
\item the surfaces meet at a positive angle along every non-seam edge;
\item at every vertex the tangent cone is an embedded PL disk.
\end{itemize}
Then a realization exists.
\end{proposition}
\begin{proof}[Proof sketch]
\begin{itemize}
\item Take $|K|$ to be the re-trimmed surface with its cells, and $h_0$ the identity.
\item By compactness, $h_0(\bar\sigma)$ has a positive margin $m$ inside the open zone $\Z(\sigma)$.
\item $h_0$ is bi-Lipschitz: within a cell by regularity, across an edge by the positive crease angle, and at a
vertex by the tangent-cone condition.
\item A secant approximation on a fine shape-regular curved triangulation is injective once its Lipschitz
distance to $h_0$ is below the bi-Lipschitz constant, and it stays in the zones once its distance is below $m$.
\end{itemize}
\end{proof}

\begin{remark}[Checking the hypotheses]
\begin{itemize}
\item The re-trimmed model can be computed by extending the face surfaces and rebuilding the solid
(OpenCascade's \texttt{BOPAlgo\_MakerVolume}). The hypotheses are then checked on samples
(\texttt{compute/reintersect\_curved.py}), except the vertex part of the transversality hypothesis, which is not
checked. That matters for the cylinder and the pipe, whose seam vertices lie on only two surfaces. The PL step is
not carried out.
\item For the planar case, the hypotheses of Proposition~\ref{prop:B} for the sewn box and the sewn plate were also
re-checked in exact rational arithmetic, including embeddedness.
\item Results are in Section~\ref{sec:experiments}.
\end{itemize}
\end{remark}
